\chapter{Estudos de Caso e Resultados Experimentais}
\label{cap:estudo_de_caso}

% TODO(P4): Detalhar a implementação prática, os estudos de caso reais e a discussão crítica dos benchmarks:
%
% 1. Estudos de Caso Práticos:
%    - Ferramenta de Diff Textual (diff_tool.py): aplicação direta de LCS considerando linhas como tokens,
%      comportamento em arquivos reais de código/texto em data/.
%    - Bioinformática e Alinhamento de DNA (dna_alignment.py): distância de edição de Levenshtein e
%      alinhamento global com inserção de gaps ('-') sobre arquivos FASTA.
%
% 2. Metodologia Experimental:
%    - Instrumentação de tempo (time.perf_counter) e memória (tracemalloc).
%    - Protocolo com mediana de 5 repetições, semente aleatória fixa e corte de 10s para algoritmos exponenciais.
%    - Metadados do ambiente de execução (hardware, CPU, sistema operacional, versões do Python e g++).
%
% 3. Análise dos Resultados Gráficos (inserir figuras geradas por plot_results.py em relatorio/figuras/):
%    - Experimento E1: Fibonacci (tempo e chamadas vs. n em escala logarítmica).
%    - Experimento E2: Corte de Hastes (\Theta(2^n) vs. \Theta(n^2)).
%    - Experimento E3: LCS - consumo de memória (tabela completa vs. duas linhas).
%    - Experimento E4: Gargalo de Linguagem - Python vs. C++17 compilado com -O3.
%    - Experimento E5: Limite de Recursão do Python (sys.getrecursionlimit() na memoização).
%    - Experimento E6: Taxa de falha empírica da heurística gulosa no corte de hastes.
%    - Experimento E7: Cadeia de Matrizes - tempo da PD vs. crescimento dos números de Catalan.
%
% 4. Discussão Crítica de Gargalos:
%    - Overhead do interpretador Python frente ao C++.
%    - Impacto da localidade espacial de cache (acesso a matrizes por linhas vs. colunas).
%    - Consumo excessivo de memória na matriz 2D em strings longas.

% [Texto e figuras a cargo de P4 - ver docs/tarefas/P4_estudo_de_caso_e_benchmarks.md]
